\documentclass[12pt]{article}
\usepackage[margin=0.9in]{geometry}
\usepackage[T1]{fontenc}
\usepackage[utf8]{inputenc}
\usepackage{hyperref}
\usepackage{xcolor}
\usepackage{enumitem}
\usepackage{parskip}
\usepackage{longtable}
\usepackage{array}
\usepackage{setspace}
\onehalfspacing

\definecolor{navy}{HTML}{003087}
\definecolor{brick}{HTML}{B22222}
\hypersetup{colorlinks=true,linkcolor=navy,urlcolor=brick}
\setlist[itemize]{leftmargin=*,itemsep=2pt,topsep=2pt}
\setlist[enumerate]{leftmargin=*,itemsep=2pt,topsep=2pt}

\title{EDU 400A Comprehensive Portfolio Rough Draft\\\large Proficiency 2: Learner Diversity \& Equity}
\author{Piter Garcia}
\date{May 2026}

\begin{document}
\maketitle

\section*{Submission note}

This document is submitted for the EDU 400A rough-draft requirement. The selected portfolio paper is \textbf{Proficiency 2: Learner Diversity and Equity}. The personal statement section is included only as brief context for the selected proficiency draft because the assignment allows either a personal statement or one proficiency paper, and this submission focuses on one proficiency paper.

\section*{Rubric, certification area, and public repository}

Rubric and certification-area requirements folder: \href{https://drive.google.com/drive/folders/1AHOAh-MUIWg-2xVZXCz6hDzydanlecPe}{Google Drive rubric folder}.

Public GitHub/Overleaf repository: \href{https://github.com/pzg8794/ED400-Comprehensive_Teaching_Portfolio}{ED400-Comprehensive\_Teaching\_Portfolio}.

Certification area: \textbf{K--12 Computer Science}. The current Drive folder contains Warner School comprehensive portfolio rubrics by certification area. Because the visible folder set does not include a separate K--12 Computer Science rubric file, this draft uses the shared Warner portfolio structure and the \textbf{Childhood Education / Elementary} rubric as the closest available working rubric for the elementary clinical evidence used here. The evidence base centers Grade 4 TeachingPlacement, Pine Brook Elementary, Minecraft Education, elementary computer science instruction, ED452B unit planning, ED400A reflection, and EDU442 identity/equity work. If a K--12 Computer Science-specific rubric is provided later, the draft should be checked against that rubric before final portfolio submission.

\tableofcontents
\newpage

\section{Portfolio structure and rough-draft focus}

The full comprehensive portfolio includes a personal statement and ten proficiency narratives. This report contains:

\begin{itemize}
  \item a short personal statement context section; and
  \item the selected EDU 400A rough draft: \textbf{Proficiency 2: Learner Diversity and Equity}.
\end{itemize}

The full future portfolio will include eleven papers: Personal Statement; Proficiency 1: Learner Development; Proficiency 2: Learner Diversity and Equity; Proficiency 3: Learning Environments; Proficiency 4: Content Knowledge and Application; Proficiency 5: Assessment of Learning; Proficiency 6: Planning for Instruction; Proficiency 7: Pedagogy; Proficiency 8: Communication and Collaboration; Proficiency 9: Reflection and Professional Growth; and Proficiency 10: Advocacy and Ethical Leadership.

\section{Artifact reference hub}

\begin{longtable}{p{0.11\textwidth}p{0.28\textwidth}p{0.31\textwidth}p{0.22\textwidth}}
\textbf{ID} & \textbf{Artifact} & \textbf{Portfolio use} & \textbf{Public link}\\ \hline
P00 & Portfolio report source & Overleaf/GitHub source for this PDF & \href{https://github.com/pzg8794/ED400-Comprehensive_Teaching_Portfolio/blob/main/reports/portfolio.tex}{portfolio.tex}\\
P02 & Proficiency 2 rough draft & EDU 400A selected rough draft & \href{https://github.com/pzg8794/ED400-Comprehensive_Teaching_Portfolio/blob/main/drafts/proficiency-02-learner-diversity-equity.md}{Markdown draft}\\
A01 & Dear Future-Teacher Me public copy & Personal Statement context; Proficiency 2, 9, and 10 evidence & \href{https://github.com/pzg8794/ED400-Comprehensive_Teaching_Portfolio/blob/main/artifacts/public-copies/dear-future-teacher-me.md}{public copy}\\
A02 & Grade 4 IGNITE Minecraft evidence summary & Proficiency 2, 3, 5, 6, 7, and 8 clinical evidence & \href{https://github.com/pzg8794/ED400-Comprehensive_Teaching_Portfolio/blob/main/artifacts/public-copies/minecraft-grade4-ignite-evidence-summary.md}{public copy}\\
A03 & Trauma-informed reflection summary & Personal Statement context; Proficiency 2, 3, 9, and 10 evidence & \href{https://github.com/pzg8794/ED400-Comprehensive_Teaching_Portfolio/blob/main/artifacts/public-copies/trauma-informed-reflection-summary.md}{public copy}\\
A04 & Childhood / Elementary rubric summary & Working rubric alignment until a CS-specific file is supplied & \href{https://github.com/pzg8794/ED400-Comprehensive_Teaching_Portfolio/blob/main/rubrics/childhood-education-rubric-summary.md}{rubric summary}\\
A05 & Rubric Drive folder manifest & Shows all rubric files found in the Drive folder & \href{https://github.com/pzg8794/ED400-Comprehensive_Teaching_Portfolio/blob/main/rubrics/drive-folder-manifest.md}{manifest}\\
\end{longtable}

Because this is a public repository, private ED400A and TeachingPlacement originals are treated as provenance. Public-safe copies and summaries are used for submission. Raw student transcripts, identifiable student work, and school-only records should not be placed in this public repository.

\newpage
\section{Personal Statement - Context Draft}

I am becoming a teacher because I understand school as both a place of possibility and a place where students can be misunderstood. My commitments come from the intersection of computer science, inclusion, culturally and linguistically sustaining teaching, neurodiversity-affirming practice, and a belief that classroom systems can either widen or close opportunity gaps.

In ED400A, I wrote to my future-teacher self that I was excited to teach in community and to build culturally and linguistically sustaining computer science lessons, trauma-informed and neurodiversity-affirming routines, and technology/AI use that lowers barriers instead of raising them. That line still names the teacher I am trying to become. I do not want technology to function as another sorting machine. I want it to become a tool students can use to think, create, collaborate, and be seen differently.

Good teaching is rigorous, humane, and responsive. It does not treat care and academic challenge as opposites. It asks what students already know, how they make meaning, what barriers are built into the task, and what forms of evidence will show their thinking fairly. For culturally and linguistically diverse and neurodivergent learners especially, good teaching requires teachers to question whether a student is actually struggling with the learning target or with the way the task was designed.

My developing conception of teaching is shaped by the idea of teacher agency as algorithmic justice. Classrooms run on systems: routines, rubrics, technologies, pacing expectations, response formats, and behavior norms. A teacher cannot control every system, but a teacher can inspect, interrupt, and redesign classroom routines so more students have access to dignity and meaningful learning.

My Grade 4 TeachingPlacement work with Minecraft Education made this concrete. Minecraft created excitement and collaboration, but it also exposed access barriers: sign-in procedures, account lockout risk, world-entry steps, reading directions from NPCs, keyboard/mouse navigation, peer coordination, and frustration after errors. Those barriers taught me that engagement is not the same as equity.

ED452B helped me turn those classroom observations into an instructional unit and evidence system. In the Minecraft Coding FUNdamentals work, I planned projected modeling, read-aloud directions, Driver/Navigator roles, planning sheets, oral explanation options, and debugging routines. These were not extra supports for a few students; they were classwide design choices that made rigorous learning more reachable.

I want to hold onto four commitments: every student belongs; relationship comes before content; rigor and care are not opposites; and technology/AI must be transparent, inclusive, and just.

\newpage
\section{Proficiency 2: Learner Diversity \& Equity - Rough Draft}

\subsection{Rubric focus}

The working rubric defines Proficiency 2 as evidence that the candidate draws from theory, research, and clinical practice to demonstrate commitments to educational equity and justice. It asks the candidate to show an understanding of how family and community backgrounds, race and ethnicity, linguistic and cultural resources, gender and sexual orientation, socioeconomic status, exceptional needs and abilities, and other personal and educational experiences affect students' ways of learning. It also asks the candidate to understand how social structures and educational institutions affect diverse learners' opportunities and experiences.

For this rough draft, I address the three rubric criteria directly: 2.1 Understanding of learner diversity and equity; 2.2 Commitment to cultivating learner diversity and equity; and 2.3 Incorporation of learner diversity and equity into inclusive, anti-oppressive learning experiences.

\subsection{Opening claim}

My teaching practice is developing around a central belief: learner diversity is not an obstacle that appears after instruction is planned; it is the starting condition for responsible instructional design. In my coursework and clinical practice, I am learning that equity means more than treating students kindly or offering help after they struggle. Equity requires designing learning experiences that recognize how students' race, language, culture, gender, disability, socioeconomic context, family background, prior educational experiences, and institutional conditions shape what they can access and how their knowledge becomes visible.

I selected Proficiency 2 because my strongest current evidence comes from work where I had to confront access and equity in concrete ways: the ED452B Minecraft Education innovative unit, its assessment map and rubric, de-identified TeachingPlacement evidence from Grade 4 IGNITE Minecraft lessons, and identity/equity reflections from ED400A and EDU442.

\subsection{Exhibit cluster}

\begin{longtable}{p{0.29\textwidth}p{0.29\textwidth}p{0.34\textwidth}}
\textbf{Exhibit} & \textbf{Source} & \textbf{Rubric connection}\\ \hline
Minecraft Education Coding FUNdamentals unit & ED452B / TeachingPlacement & Required exhibit: innovative unit, materials, and assessment rubric; supports 2.1, 2.2, and 2.3\\
Assessment map and evidence system & ED452B & Required/recommended exhibit: assessment rubric and student assessment data system; supports 2.1 and 2.3\\
De-identified Grade 4 IGNITE placement evidence & TeachingPlacement & Recommended exhibit: clinical lesson/unit evidence, field notes, student work, and assessment data; supports 2.3\\
ED400A and EDU442 identity/equity reflections & ED400A / EDU442 & Recommended exhibit: key papers, reflections, and coursework artifacts; supports 2.1 and 2.2\\
\end{longtable}

\subsection{2.1 Understanding of learner diversity and equity}

The rubric expects me to understand how learner diversity includes race and ethnicity, language use, gender and sexual orientation, socioeconomic status, disability status, personal and cultural characteristics, and the social structures and institutions that surround students. My current understanding is that these categories do not sit outside the classroom waiting to be mentioned during an equity lesson. They shape the ordinary conditions of learning: who feels seen, who has background knowledge treated as valuable, who understands the language of directions, who has enough time to process, who can use the required tool, who is labeled as capable, and who is misread as disengaged.

The Minecraft Education unit helped me see this clearly. On the surface, Minecraft looks like a high-interest, student-friendly digital tool. In practice, it can either expand access or reproduce inequity depending on how the lesson is designed. A student may understand sequencing and debugging but struggle to complete the task because the directions are text-heavy, the interface is unfamiliar, the keyboard/mouse demands are high, or the pace assumes executive-function skills that not every student has developed in the same way. If I only look at the final product, I can confuse access barriers with conceptual misunderstandings.

That distinction matters for learner diversity and equity. Culturally and linguistically diverse learners may need oral directions, visual models, repeated language, or chances to explain thinking through modes other than formal written English. Neurodivergent learners may need visible routines, reduced working-memory load, structured partner roles, or explicit debugging steps. Students with different socioeconomic and digital access histories may need time to learn the platform before the platform is used to judge their content understanding.

\subsection{2.2 Commitment to cultivating learner diversity and equity}

The rubric asks not only whether I understand equity, but whether I show a commitment to cultivating it. That commitment has developed across my teacher education program through coursework, reflection, clinical practice, and my own lived experience as a neurodivergent, culturally and linguistically complex learner. My ED400A and EDU442 work repeatedly returns to the same professional stance: difference is not deficiency, and school systems often misread difference before they support it.

In ED400A, my reflections around teacher agency, trauma-informed teaching, and advocacy pushed me to think about equity as a professional responsibility rather than a personal preference. I learned to ask how schools interpret behavior, language, disability, and family context. A student who is quiet, off task, impulsive, or frustrated is not automatically refusing to learn. They may be responding to overload, unclear routines, language demands, disability-related needs, cultural mismatch, trauma, or an environment that was not designed with them in mind.

The Minecraft unit reflects this commitment by treating access supports as normal parts of rigorous learning. Projected click paths, read-aloud of NPC directions, planning sheets, Driver/Navigator roles, sentence frames, oral explanation options, and a Stop--Read--Predict--Run--Debug protocol are not add-ons for a few students. They are classwide structures that protect dignity and allow more students to participate without being publicly marked as the problem.

\subsection{2.3 Inclusive, anti-oppressive learning experiences}

The third rubric criterion asks whether I can incorporate learner diversity and equity into inclusive, anti-oppressive learning experiences. This is where the portfolio has to move from belief to practice. The ED452B Minecraft unit and TeachingPlacement evidence show my attempt to build a learning experience where students could access computational thinking through multiple routes.

The unit was designed around elementary computer science concepts: algorithms, sequencing, debugging, efficiency, and reflection. To make those concepts accessible, I planned multiple representations and response modes. Students could see modeled click paths, hear directions read aloud, use planning sheets before coding, work with a partner, explain their thinking orally, revise code after testing, and use debugging as an expected part of learning rather than a sign of failure.

The assessment map strengthened the inclusive design because it refused to reduce learning to task completion. It separated several questions: Did the student access the Minecraft world? Did the student understand the goal? Did the student plan a sequence? Did the student test and debug? Could the student explain the computational idea? This matters because a student who fails to finish the world may still understand the concept, and a student who finishes quickly may still need to explain their reasoning.

The TeachingPlacement evidence helped me see that inclusive design has to stay responsive. In the Grade 4 IGNITE Minecraft lessons, students encountered predictable barriers: login and navigation issues, NPC direction reading, rushing before planning, frustration after errors, and uneven partner participation. These patterns pushed me to adjust support in real time. I began to see student struggle as information about the system, not simply as information about the child.

\subsection{Growth and conclusion}

My biggest growth is that I now see access as part of content instruction. Earlier in my preparation, I might have treated supports as accommodations or modifications that happen after a lesson is already designed. Now I see them as part of the architecture of rigorous teaching. A visual model, a planning sheet, an oral explanation option, or a classwide debugging routine does not lower the learning target. It makes the target reachable and visible.

The artifacts in this rough draft show that I am developing a concrete, evidence-based approach to learner diversity and equity. I understand equity as attention to both individual learner differences and the social/institutional structures that shape opportunity. I am committed to designing learning experiences where students' identities, languages, abilities, interests, and histories are treated as resources. I am learning to enact that commitment through inclusive, anti-oppressive practices: classwide access supports, multiple ways to demonstrate learning, careful assessment interpretation, de-identified student evidence, and reflection that inspects the system before blaming the learner.

For EDU 400A, this draft is intentionally rough. The next revision should add more direct theory/research citations, replace any remaining summaries with final public PDF artifacts where appropriate, and re-check the language against a K--12 Computer Science rubric if one becomes available.

\end{document}
